\documentclass[fleqn]{article}

\usepackage{aima-slides}
\usepackage[utf8]{inputenc}
\usepackage[T5]{fontenc}
\usepackage{lmodern}

\begin{document}

\begin{huge}
\titleslide{Rational decisions}{Chapter 16}

%%%%%%%%%%%% Slide %%%%%%%%%%%%%%%%%%%%%%%%%%%%%%%%%%%%%%%%%%%%%%%%%%%
\heading{Outline}

\blob Rational preferences

\blob Utilities

\blob Money

\blob Multiattribute utilities

\blob Decision networks

\blob Value of information



%%%%%%%%%%%% Slide %%%%%%%%%%%%%%%%%%%%%%%%%%%%%%%%%%%%%%%%%%%%%%%%%%%
\heading{Preferences}

An agent chooses among \u{prizes} ($A$, $B$, etc.) and 
\u{lotteries}, i.e., situations with uncertain prizes

\begin{tabular}{lr}
\hbox{\begin{minipage}[b]{0.6\textwidth}
\vspace*{0.3in}

Lottery $L = [p,A;\ (1-p),B]$

\vspace*{0.3in}

\end{minipage}}
&
\epsfxsize=0.3\textwidth
\epsffile{\file{figures}{lottery.ps}}
\end{tabular}

Notation:\al
$A \pref B$ \qquad $A$ preferred to $B$\al
$A \indiff B$ \qquad indifference between $A$ and $B$\al
$A \prefeq B$ \qquad $B$ not preferred to $A$



%%%%%%%%%%%% Slide %%%%%%%%%%%%%%%%%%%%%%%%%%%%%%%%%%%%%%%%%%%%%%%%%%%
\heading{Rational preferences}

Idea: preferences of a rational agent must obey constraints.\\
Rational preferences $\implies$ \nl
   behavior describable as maximization of expected utility

Constraints:\al
\underline{Orderability}\nl
$(A \pref B) \lor (B \pref A) \lor (A \indiff B)$\al
\underline{Transitivity}\nl
$(A \pref B) \land (B \pref C) \implies (A \pref C)$\al
\underline{Continuity}\nl
$A \pref B \pref C \implies \Exi{p} [p,A;\ 1-p,C] \indiff B$\al
\underline{Substitutability}\nl
$A \indiff B \implies [p,A;\ 1-p,C] \indiff [p,B; 1-p,C]$\al
\underline{Monotonicity}\nl
$A \pref B \implies (p \geq q \lequiv [p,A;\ 1-p,B] \prefeq [q,A;\ 1-q,B])$




%%%%%%%%%%%% Slide %%%%%%%%%%%%%%%%%%%%%%%%%%%%%%%%%%%%%%%%%%%%%%%%%%%
\heading{Rational preferences contd.}

Violating the constraints leads to self-evident irrationality

For example: an agent with intransitive preferences
can be induced to give away all its money

\begin{tabular}{lr}
\hbox{\begin{minipage}[b]{0.5\textwidth}

If $B \pref C$, then an agent who has $C$
would pay (say) 1 cent to get $B$

\vspace*{0.2in}

If $A \pref B$, then an agent who has $B$
would pay (say) 1 cent to get $A$

\vspace*{0.2in}

If $C \pref A$, then an agent who has $A$
would pay (say) 1 cent to get $C$


\end{minipage}}
&
\epsfxsize=0.3\textwidth
\ \qquad\qquad\epsffile{\file{figures}{cash-machine.ps}}
\end{tabular}




%%%%%%%%%%%% Slide %%%%%%%%%%%%%%%%%%%%%%%%%%%%%%%%%%%%%%%%%%%%%%%%%%%
\heading{Maximizing expected utility}

\u{Theorem} (Ramsey, 1931; von Neumann and Morgenstern, 1944):\\
Given preferences satisfying the constraints\\
there exists a real-valued function $U$ such that\nl
    $U(A) \geq U(B)\ \lequiv \ A\prefeq B$\nl
    $U([p_1,S_1;\ \ldots\ ;\ p_n,S_n]) = \mysum_i\ p_i U(S_i)$

\u{MEU principle}:\al
Choose the action that maximizes expected utility

Note: an agent can be entirely rational (consistent with MEU)\\
without ever representing or manipulating utilities and probabilities

E.g., a lookup table for perfect tictactoe

%%%%%%%%%%%% Slide %%%%%%%%%%%%%%%%%%%%%%%%%%%%%%%%%%%%%%%%%%%%%%%%%%%
\heading{Utilities}

Utilities map states to real numbers. Which numbers?

Standard approach to assessment of human utilities:\al
  compare a given state $A$ to a \u{standard lottery} $L_p$ that has\nl
    ``best possible prize'' $\ubest$ with probability $p$\nl
    ``worst possible catastrophe'' $\uworst$ with probability $(1-p)$\al
  adjust lottery probability $p$ until $A \indiff L_p$

\vspace*{0.4in}

\epsfxsize=0.85\textwidth
\fig{\file{figures}{micromort.ps}}

%%%%%%%%%%%% Slide %%%%%%%%%%%%%%%%%%%%%%%%%%%%%%%%%%%%%%%%%%%%%%%%%%%
\heading{Utility scales}

\u{Normalized utilities}: $\ubest = 1.0$, $\uworst = 0.0$

\u{Micromorts}: one-millionth chance of death\al
  useful for Russian roulette, paying to reduce product risks, etc.

\u{QALYs}: quality-adjusted life years\al
  useful for medical decisions involving substantial risk

Note: behavior is \u{invariant} w.r.t. +ve linear transformation
\[
  U'(x) = k_1 U(x) + k_2 \quad\mbox{where } k_1 > 0
\]
With deterministic prizes only (no lottery choices), only\\
\u{ordinal utility} can be determined, i.e., total order on prizes

%%%%%%%%%%%% Slide %%%%%%%%%%%%%%%%%%%%%%%%%%%%%%%%%%%%%%%%%%%%%%%%%%%
\heading{Money}

Money does \u{not} behave as a utility function

Given a lottery $L$ with expected monetary value $EMV(L)$,\\
usually $U(L) < U(EMV(L))$, i.e., people are \u{risk-averse}

Utility curve: for what probability $p$ am I indifferent between\
a fixed prize $x$ and a lottery $[p,\$M;\ (1-p),\$0]$ for large $M$?

Typical empirical data, extrapolated with \u{risk-prone} behavior:

\vspace*{0.2in}

\epsfxsize=0.55\textwidth
\fig{\file{figures}{beard-utility.ps}}



%%%%%%%%%%%% Slide %%%%%%%%%%%%%%%%%%%%%%%%%%%%%%%%%%%%%%%%%%%%%%%%%%%
\heading{Student group utility}

For each $x$, adjust $p$ until half the class votes for lottery (M=10,000)

\vspace*{0.2in}

\epsfxsize=1.05\textwidth
\fig{\file{figures}{student-utility.ps}}



%%%%%%%%%%%% Slide %%%%%%%%%%%%%%%%%%%%%%%%%%%%%%%%%%%%%%%%%%%%%%%%%%%
\heading{Decision networks}

Add \u{action nodes} and \u{utility} nodes to belief networks\\
to enable rational decision making

\vspace*{0.2in}

\epsfxsize=0.48\textwidth
\fig{\file{figures}{airport-id.ps}}

Algorithm:\al
  For each value of action node\nl
    compute expected value of utility node given action, evidence\al
  Return MEU action


%%%%%%%%%%%% Slide %%%%%%%%%%%%%%%%%%%%%%%%%%%%%%%%%%%%%%%%%%%%%%%%%%%
\heading{Multiattribute utility}

How can we handle utility functions of many variables $X_1\ldots X_n$?\\
E.g., what is $U(Deaths,Noise,Cost)$?

How can complex utility functions be assessed from \\
preference behaviour?

Idea 1: identify conditions under which decisions can be made without
complete identification of $U(x_1,\ldots,x_n)$

Idea 2: identify various types of \u{independence} in preferences\\
and derive consequent canonical forms for $U(x_1,\ldots,x_n)$

%%%%%%%%%%%% Slide %%%%%%%%%%%%%%%%%%%%%%%%%%%%%%%%%%%%%%%%%%%%%%%%%%%
\heading{Strict dominance}

Typically define attributes such that $U$ is \u{monotonic} in each

\u{Strict dominance}: choice $B$ strictly dominates choice $A$ iff\nl
$\All{i} X_i(B) \geq X_i(A)$ \quad (and hence $U(B) \geq U(A)$)

\vspace*{0.2in}

\epsfxsize=0.8\textwidth
\fig{\file{figures}{strict-dominance.ps}}

Strict dominance seldom holds in practice


%%%%%%%%%%%% Slide %%%%%%%%%%%%%%%%%%%%%%%%%%%%%%%%%%%%%%%%%%%%%%%%%%%
\heading{Stochastic dominance}

\twograph{\file{graphs}{dominance-density.ps}}{\file{graphs}{dominance-cumulative.ps}}

Distribution $p_1$ \u{stochastically dominates} distribution $p_2$ iff\nl
  $\displaystyle\All{t} \int_{-\infty}^t p_1(x)dx \leq \int_{-\infty}^t p_2(t)dt$

If $U$ is monotonic in $x$, then $A_1$ with outcome distribution $p_1$\\
stochastically dominates $A_2$ with outcome distribution $p_2$:\nl
  $\displaystyle\int_{-\infty}^{\infty} p_1(x) U(x)dx \geq \int_{-\infty}^{\infty} p_2(x) U(x)dx $\\
Multiattribute case: stochastic dominance on all attributes $\implies$ optimal



%%%%%%%%%%%% Slide %%%%%%%%%%%%%%%%%%%%%%%%%%%%%%%%%%%%%%%%%%%%%%%%%%%
\heading{Stochastic dominance contd.}

Stochastic dominance can often be determined without\\
exact distributions using \u{qualitative} reasoning

E.q., construction cost increases with distance from city\nl
      $S_2$ is further from the city than $S_1$\al
  $\implies$ $S_1$ stochastically dominates $S_2$ on cost

E.g., injury increases with collision speed

Can annotate belief networks with stochastic dominance information:\al
  $X \qplus Y$ ($X$ positively influences $Y$) means that\al
  For every value $\mbf{z}$ of $Y$'s other parents $\mbf{Z}$\nl
    $\All{x_1,x_2} x_1 \geq x_2 \implies
      \pv(Y|x_1,\mbf{z})$ stochastically dominates $ \pv(Y|x_2,\mbf{z})$


%%%%%%%%%%%% Slide %%%%%%%%%%%%%%%%%%%%%%%%%%%%%%%%%%%%%%%%%%%%%%%%%%%
\heading{Example: car insurance}

Which arcs are positive or negative influences?

\vspace*{0.2in}

\epsfxsize=0.95\textwidth
\fig{\file{figures}{insurance-net.ps}}


%%%%%%%%%%%% Slide %%%%%%%%%%%%%%%%%%%%%%%%%%%%%%%%%%%%%%%%%%%%%%%%%%%
\heading{Preference structure: Deterministic}

$X_1$ and $X_2$ \u{preferentially independent} of $X_3$ iff\al
  preference between $\< x_1,x_2,x_3 \>$ and $\< x_1',x_2',x_3 \>$\al
  does not depend on $x_3$

E.g., $\<Noise,Cost,Safety\>$:\al
  $\<$20,000 suffer, \$4.6 billion, 0.06 deaths/mpm$\>$ vs.\al
  $\<$70,000 suffer, \$4.2 billion, 0.06 deaths/mpm$\>$

\u{Theorem} (Leontief, 1947): if every pair of attributes is P.I. of its complement,
then every subset of attributes is P.I of its complement: \u{mutual P.I.}.

\u{Theorem} (Debreu, 1960): mutual P.I. $\implies$ $\exists$ \u{additive} value function:
\[
  V(S) = \mysum_i V_i(X_i(S))
\]
Hence assess $n$ single-attribute functions; often a good approximation


%%%%%%%%%%%% Slide %%%%%%%%%%%%%%%%%%%%%%%%%%%%%%%%%%%%%%%%%%%%%%%%%%%
\heading{Preference structure: Stochastic}

Need to consider preferences over lotteries:\\
$\mbf{X}$ is \u{utility-independent} of $\mbf{Y}$ iff\al
  preferences over lotteries $\mbf{X}$ do not depend on $\mbf{y}$

Mutual U.I.: each subset is U.I of its complement\\
$\implies$ $\exists$ \u{multiplicative} utility function:\al
$U  =  k_1U_1 + k_2U_2 + k_3U_3$\nl
 + $k_1k_2U_1U_2 + k_2k_3U_2U_3 + k_3k_1U_3U_1$\nl
 + $k_1k_2k_3U_1U_2U_3$

Routine procedures and software packages for generating preference
tests to identify various canonical families of utility functions


%%%%%%%%%%%% Slide %%%%%%%%%%%%%%%%%%%%%%%%%%%%%%%%%%%%%%%%%%%%%%%%%%%
\heading{Value of information}

Idea: compute value of acquiring each possible piece of evidence\\
Can be done \u{directly from decision network}

Example: buying oil drilling rights\al
  Two blocks $A$ and $B$, exactly one has oil, worth $k$\al
  Prior probabilities 0.5 each, mutually exclusive\al
  Current price of each block is $k/2$\al
  Consultant offers accurate survey of $A$. Fair price?

Solution: compute expected value of information\al
  = expected value of best action given the information\nl
    minus expected value of best action without information\\
Survey may say ``oil in A'' or ``no oil in A'', prob. 0.5 each\al
  = [$0.5 \times {}$ value of ``buy A'' given ``oil in A''\nl
    + $0.5 \times {}$ value of ``buy B'' given ``no oil in A'']\nl
    -- 0\al
  = $(0.5 \times k/2) + (0.5 \times k/2) - 0 = k/2$



%%%%%%%%%%%% Slide %%%%%%%%%%%%%%%%%%%%%%%%%%%%%%%%%%%%%%%%%%%%%%%%%%%
\heading{General formula}

Current evidence $E$, current best action $\alpha$\\
Possible action outcomes $S_i$, potential new evidence $E_j$
\[
  EU(\alpha|E) = \max_{a} \mysum_i\ U(S_i)\;P(S_i|E,a)
\]
Suppose we knew $E_j \eq e_{jk}$, then we would choose $\alpha_{e_{jk}}$ s.t.
\[
  EU(\alpha_{e_{jk}}|E,E_j \eq e_{jk}) = \max_a \mysum_i\ U(S_i)\;P(S_i|E,a,E_j \eq e_{jk})
\]
$E_j$ is a random variable whose value is {\it currently} unknown\\
$\implies$ must compute expected gain over all possible values:
\[
VPI_{E}(E_j) = \left(\mysum_k\ P(E_j \eq e_{jk}|E)
EU(\alpha_{e_{jk}}|E,E_j \eq e_{jk})\right) - EU(\alpha|E) 
\]
(VPI = value of perfect information)


%%%%%%%%%%%% Slide %%%%%%%%%%%%%%%%%%%%%%%%%%%%%%%%%%%%%%%%%%%%%%%%%%%
\heading{Properties of VPI}

\u{Nonnegative}---in {\em expectation}, not {\em post hoc}
\[
\All{j,E} VPI_{E}(E_j)\geq 0\
\]
\u{Nonadditive}---consider, e.g., obtaining $E_j$ twice
\[
VPI_{E}(E_j,E_k) \not= VPI_{E}(E_j) + VPI_{E}(E_k)
\]
\u{Order-independent}
\[
VPI_{E}(E_j,E_k) = VPI_{E}(E_j) + VPI_{E,E_j}(E_k)
                   = VPI_{E}(E_k) + VPI_{E,E_k}(E_j) 
\]
Note: when more than one piece of evidence can be gathered,\\
maximizing VPI for each to select one is not always optimal\\
$\implies$ evidence-gathering becomes a \u{sequential} decision problem



%%%%%%%%%%%% Slide %%%%%%%%%%%%%%%%%%%%%%%%%%%%%%%%%%%%%%%%%%%%%%%%%%%
\heading{Qualitative behaviors}

a) Choice is obvious, information worth little\\
b) Choice is nonobvious, information worth a lot\\
c) Choice is nonobvious, information worth little

\vspace*{-0.2in} %% the following figure's bounding box includes whitespace

\epsfxsize=1.0\textwidth
\fig{\file{figures}{3cases.ps}}





\end{huge} 
\end{document}
